\chapter{Definition of the Bulk-Silicon Program}
\index{silicon!bulk program|textbf}

\label{ch:physical-problem}

\chapterstatus{expository synthesis of the accepted bulk-silicon physical
specification; production outcomes remain proposed work.}

This chapter defines the physical parent, pristine host, and representation
sequence for the bulk-silicon program. Dopant-specific branches are introduced
separately in Part~III.

\section{The parent description}
\index{parent model!Kohn--Sham}
\index{Kohn--Sham operator}


The program asks when a first-principles Kohn--Sham description of silicon and
substitutional dopants can be replaced by a smaller model without losing the
operator information needed for its declared use. The question is not whether
one model is universally ``correct.'' It is whether a stated reduced-model
class preserves stated features of a stated parent problem within separately
declared errors.

The accepted physical conventions are owned by
\path{specification/ksdft2Effmass.physical-specification.v1.md}. This chapter
organizes those conventions as a research narrative and introduces no new
physical branch.

The parent is a self-consistent Kohn--Sham density-functional calculation in
the Born--Oppenheimer fixed-ion framework~\cite{hohenberg1964,kohn1965}. Its
output is a one-particle Kohn--Sham operator and associated finite numerical
representations. It is not the complete interacting many-body Hamiltonian, and
its eigenvalues are not identified here with the complete many-body excitation
spectrum.

\citationtodo{high priority}{Check the distinction between Kohn--Sham
eigenvalue gaps and physical fundamental or quasiparticle gaps against
Perdew--Levy (1983) and Sham--Schluter (1983), bibliography keys
\texttt{perdewLevy1983} and \texttt{shamSchluter1983} (already present).}

For the version-one program, exchange and correlation are described with the
PBE generalized-gradient approximation~\cite{perdew1996}, and the
electron--ion interaction is represented by compatible optimized
norm-conserving Vanderbilt pseudopotentials from the approved PseudoDojo
family~\cite{hamann2013,vansetten2018}. These choices define the parent against
which later reductions are measured. In particular, disagreement between this
parent and experiment is a parent-model question, not automatically an error
introduced by tight-binding or effective-mass reduction.

\section{Pristine silicon host}
\index{silicon!pristine host}
\index{diamond structure}


The host is crystalline silicon in the diamond structure: an fcc Bravais
lattice with a two-atom basis. Primitive-cell and conventional-cell
coordinates may both be useful, but every retained representation must
identify which lattice vectors, axes, basis sites, and reciprocal-coordinate
convention it uses.

The immediate bulk pilot is scalar relativistic, non-spin-polarized, and does
not include spin--orbit coupling. Its production lattice constant is to be
obtained from a zero-pressure PBE bulk relaxation using the selected silicon
pseudopotential. Until that calculation and its convergence record exist, no
numerical lattice constant is a calculated result of this project.

The pilot targets a band-edge subspace sufficient for an orthogonal
$sp^3s^{\ast}$ reduction, including the valence and low-conduction states
needed to describe the indirect Kohn--Sham gap and electron effective masses.
The exact band count, projection choices, and Wannier windows belong to later
methodological records. They are not fixed by the phrase ``band-edge
subspace'' alone.

The declared bulk observables are the indirect Kohn--Sham gap, the
conduction-valley position, longitudinal and transverse electron effective
masses, and withheld-point band errors. These quantities serve distinct roles:
they characterize the selected parent, test numerical convergence, and later
evaluate reduced representations. Passing one role does not imply passing the
others.

\section{From parent models to reduced models}
\index{reduction hierarchy}


The program studies a sequence of representations rather than one direct leap
from density-functional theory to continuum theory:
\begin{enumerate}
    \item a converged Kohn--Sham parent problem;
    \item a retained projected or Wannier subspace;
    \item an aligned finite operator representation;
    \item a constrained lattice model, initially an orthogonal
          $sp^3s^{\ast}$ hierarchy for bulk silicon;
    \item an extracted and reduced impurity operator; and
    \item where justified, a continuum effective-mass description.
\end{enumerate}

Projection, localization, alignment, parameter fitting, truncation, and
continuum approximation are different operations. Their errors and assumptions
must therefore remain identifiable. Wannier localization alone is not a
low-rank approximation, and a good spectral fit alone does not establish
agreement of aligned operator matrix elements.

\section{Bulk research questions}
\index{silicon!bulk research questions}


The bulk program asks which finite subspace preserves the declared silicon
features, whether a prescribed tight-binding class can satisfy both spectral
and aligned-operator criteria, and which tested class is the least complex one
that does so. These are proposed questions rather than completed findings.

\section{Interpretive limits}

A finite periodic supercell is not an isolated impurity. A Kohn--Sham band gap
is not an exact quasiparticle gap. Spectral agreement is not operator equality.
Matrix subtraction is not physically meaningful before common-space and
energy alignment. Software execution and passing tests establish neither
numerical convergence nor scientific validation unless the corresponding
evidence and acceptance rules are present.

\citationtodo{medium priority}{Check the isolated-impurity limitation for a
finite periodic supercell against defect-physics literature. Candidate:
Freysoldt et al. (2014), bibliography key \texttt{freysoldt2014} (already
present).}

These limits are part of the scientific problem rather than editorial caveats:
they determine which comparisons are well posed and which conclusions the
resulting hierarchy can support.
